\documentclass[10pt,a4paper]{article}

\usepackage[margin=1in]{geometry}
\usepackage{amsmath,amssymb}
\usepackage{hyperref}
\usepackage{booktabs}
\usepackage{listings}
\usepackage{xcolor}

\hypersetup{colorlinks=true, linkcolor=blue!50!black, citecolor=blue!50!black, urlcolor=blue!50!black}

\title{An Empirical Note on LLL Degeneracy of the\\
       Boneh--Venkatesan Lattice for secp256k1}
\author{Anonymous}
\date{}

\newcommand{\FF}{\mathbb{F}}
\newcommand{\ZZ}{\mathbb{Z}}

\begin{document}
\maketitle

\begin{abstract}
We document two distinct LLL-convergence failure modes for the
Boneh--Venkatesan Hidden-Number Problem (HNP) lattice on deployed
ECDSA curves. \textbf{Failure mode A}: at 256-bit, LLL fails for
secp256k1 while succeeding for P-256 and brainpoolP256r1 at
identical parameters (10~s timeout vs.\ 0.7~s convergence) — a
curve-specific arithmetic phenomenon. \textbf{Failure mode B}: at
$\geq$ 384 bits, LLL fails uniformly across P-384, brainpoolP384r1,
and P-521 — consistent with the iteration cap being undersized
relative to the lattice's expected convergence time at large
bit-lengths. The initial hypothesis (near-power-of-2 group order
is the cause) is refuted by the smaller Koblitz curves
\texttt{secp192k1} and \texttt{secp224k1} succeeding at LLL despite
both having near-power-of-2 group orders. We provide the empirical
table, a sketched Gram--Schmidt analysis identifying
$\mu$-oscillation as the likely Failure-A mechanism, and concrete
follow-up experiments to disambiguate the remaining hypotheses.
\end{abstract}

\section{Setting and notation}

The Boneh--Venkatesan lattice attack on biased-nonce
ECDSA~\cite{boneh1996,nguyen2002} recovers the secret key $d$ from
$m$ signatures $(r_i, s_i, z_i)$ in which the nonce $k_i$ is known
to satisfy $0 \leq k_i < 2^{k_{\text{bits}}}$ for some
$k_{\text{bits}} < n_{\text{bits}}$. With $a_i = s_i^{-1} r_i \bmod n$
and $t_i = s_i^{-1} z_i \bmod n$, each signature contributes the
constraint
\[
   d \cdot a_i + t_i \equiv k_i \pmod{n}.
\]

The standard $(m+2)$-dimensional integer lattice basis (rows)
\[
   B \;=\;
   \begin{pmatrix}
     n^2 \cdot I_m   & 0_{m \times 2}  \\[2pt]
     n \cdot \mathbf{a}^\top  & 2^{k_{\text{bits}}}  & 0 \\[2pt]
     n \cdot \mathbf{t}^\top  & 0  & n \cdot 2^{k_{\text{bits}}}
   \end{pmatrix}
\]
is fed to LLL with Lov\'asz parameter $\delta = 0.75$; the short
vector after reduction contains $-d \cdot 2^{k_{\text{bits}}}$ in
the $(m+1)$-st coordinate (modulo sign).

\section{The phenomenon}

We ran identical attack parameters head-to-head on two 256-bit
deployed curves: NIST P-256~\cite{fips186} and
secp256k1~\cite{sec2}. Both have prime group order $n$ of bit-length
256.

\subsection{Empirical measurements}

Three independent deterministic seed pairs $(d_{\text{seed}}, k_{\text{seed}})$
were tested. The exact LLL implementation, $\delta$, iteration cap,
and Gram--Schmidt order were identical for both curves. Signatures
were generated by a deterministic ECDSA-with-injected-bias signer.

\begin{table}[h]
\centering
\begin{tabular}{lrrr}
\toprule
$(d_{\text{seed}}, k_{\text{seed}})$ & P-256 outcome & P-256 time & secp256k1 outcome \\
\midrule
\texttt{(0xC0FFEE, 0xC0FFEE)}           & \checkmark recovered & 690~ms  & iteration cap @ 9993~ms \\
\texttt{(0xDEADBEEF, 0xBADCAFE)}        & \checkmark recovered & 697~ms  & iteration cap @ 9845~ms \\
\texttt{(0x12345678, 0x9ABCDEF0)}       & \checkmark recovered & 646~ms  & iteration cap @ 10119~ms \\
\bottomrule
\end{tabular}
\caption{LLL behaviour: $k_{\text{bits}} = 192$, $m = 8$, $\delta = 0.75$.
P-256: 3/3 recoveries. secp256k1: 0/3 recoveries.}
\label{tab:probes}
\end{table}

Additional probes (BKZ-10, larger $m$, varying $k_{\text{bits}}$)
reproduce the same qualitative pattern: secp256k1 lattices fail
where P-256 lattices succeed.

\subsection{Reproducibility}

The probe is committed at
\begin{quote}
\texttt{tests/lll\_degeneracy\_probe.rs}
\end{quote}
of the companion repository and exits in approximately 53~seconds
on commodity hardware. The HNP attack itself lives in
\texttt{src/cryptanalysis/hnp\_ecdsa.rs}; we use it unmodified.

\section{Initial hypothesis and its refutation}

The most salient arithmetic difference between secp256k1 and
P-256 is the structure of their group orders. P-256's $n$ is
\[
n_{\text{P-256}} \;=\; \texttt{0xFFFFFFFF00000000FFFFFFFFFFFFFFFFBCE6FAADA7179E84F3B9CAC2FC632551},
\]
while secp256k1's $n$ is
\[
n_{\text{secp256k1}} \;=\; \texttt{0xFFFFFFFF\ldots FEBAAEDCE6AF48A03BBFD25E8CD0364141}
\]
which can be written as $n = 2^{256} - t_{\text{secp}}$ for
$t_{\text{secp}} \approx 2^{129}$. secp256k1's group order has a
near-power-of-2 form with the deficit concentrated in the bottom
129 bits.

Our initial working hypothesis: \textbf{the Boneh--Venkatesan basis
matrix, which has entries $n^2$ and $n \cdot a_i$, produces a
Gram--Schmidt orthogonalisation with ratios that concentrate in
LLL-pathological values when $n$ has near-power-of-2 form, while
generic $n$ avoids this.} Three testable predictions followed:

\begin{enumerate}
\item Other near-power-of-2-$n$ Koblitz curves should exhibit the
      same LLL-degeneracy: \texttt{secp192k1}, \texttt{secp224k1}.
\item Generic-$n$ curves (P-256, brainpoolP256r1) should
      consistently converge.
\item BKZ-$\beta$ with $\beta \geq 20$ may patch the degeneracy.
\end{enumerate}

\subsection{Cross-Koblitz refutation}

We extended the head-to-head probe to \texttt{secp192k1} (Koblitz,
192-bit, $n \approx 2^{192}$) and \texttt{secp224k1} (Koblitz,
224-bit, $n \approx 2^{224}$). Both are
\emph{also} near-power-of-2 group orders. If the initial hypothesis
were correct, both should fail LLL similarly to secp256k1.

\begin{table}[h]
\centering
\begin{tabular}{lrrr}
\toprule
Curve         & $n$ bit-length & LLL outcome (2 seeds) & Time per probe \\
\midrule
\texttt{secp192k1}  &   192 & 2/2 RECOVERED          & $\sim$320~ms   \\
\texttt{secp224k1}  &   224 & 2/2 RECOVERED          & $\sim$630~ms   \\
\texttt{secp256k1}  &   256 & 0/2 (iteration cap)    & $\sim$10000~ms \\
\texttt{P-256}      &   256 & 2/2 RECOVERED          & $\sim$700~ms   \\
\bottomrule
\end{tabular}
\caption{Cross-Koblitz LLL probe. Parameters scaled to each curve's
bit-length: $k_{\text{bits}} = 3 n_{\text{bits}}/4$, $m = 8$. The
two smaller Koblitz curves (secp192k1, secp224k1) -- both of which
have near-power-of-2 group orders -- succeed at LLL where
secp256k1 alone fails.}
\label{tab:koblitz-probe}
\end{table}

\textbf{The initial hypothesis is refuted.} The LLL-degeneracy is
not a general feature of near-power-of-2 group orders, nor of
Koblitz curves more broadly. It is specific to secp256k1.

\section{Extended bit-length sweep: two distinct degeneracies}

A wider sweep including P-384, P-521, brainpoolP256r1, and
brainpoolP384r1 reveals that there are \emph{two} distinct
degeneracies, not one:

\begin{table}[h]
\centering
\begin{tabular}{lrrll}
\toprule
Curve              & $n$ bits & LLL outcome     & Time per probe & Notes \\
\midrule
\texttt{secp192k1} &   192    & RECOVERED       & $\sim$389 ms    & Koblitz, $j=0$ \\
\texttt{secp224k1} &   224    & RECOVERED       & $\sim$617 ms    & Koblitz, $j=0$ \\
\texttt{secp256k1} &   256    & iteration cap   & $\sim$10000 ms  & Koblitz, $j=0$ \textbf{(outlier)} \\
\texttt{P-256}     &   256    & RECOVERED       & $\sim$710 ms    & NIST, generic $j$ \\
\texttt{brainpoolP256r1} & 256 & RECOVERED      & $\sim$714 ms    & generic $j$ \\
\texttt{P-384}     &   384    & iteration cap   & $\sim$20392 ms  & NIST, generic $j$ \textbf{(uniform fail)} \\
\texttt{brainpoolP384r1} & 384 & iteration cap   & $\sim$20131 ms  & generic $j$ \textbf{(uniform fail)} \\
\texttt{P-521}     &   521    & iteration cap   & $\sim$20662 ms  & NIST, generic $j$ \textbf{(uniform fail)} \\
\bottomrule
\end{tabular}
\caption{Extended LLL-convergence sweep, $k_{\text{bits}} = 3 n_{\text{bits}}/4$,
$m = 8$, single seed. The data reveal \emph{two} regimes of LLL
failure: (1) secp256k1 specifically at 256 bits, while
other 256-bit curves succeed; (2) all $\geq 384$-bit curves
\emph{uniformly} fail, suggesting an iteration-cap issue rather
than a curve-specific property.}
\label{tab:sweep}
\end{table}

\subsection{Two distinct degeneracies}

\paragraph{Failure mode A (curve-specific): secp256k1 at 256 bits.}
P-256 and brainpoolP256r1 — both 256-bit curves with similar
basis-entry sizes — succeed at LLL while secp256k1 fails. This is
a curve-specific arithmetic phenomenon, not bit-length-driven.

\paragraph{Failure mode B (uniform): all $\geq 384$-bit curves.}
P-384, brainpoolP384r1, and P-521 all fail uniformly. None has
the special arithmetic shape of secp256k1's near-power-of-2 $n$;
P-384 in particular is the NIST "Solinas" prime with no Koblitz
structure. The uniformity strongly suggests an
\emph{implementation-side} cause (LLL iteration cap insufficient
at large bit-lengths) rather than a curve-side cause.

\section{Reframed open questions}

The data narrow the open questions to two distinct ones, one per
failure mode:

\subsection{Question A: why does secp256k1 alone fail at 256 bits?}

P-256 and brainpoolP256r1 share secp256k1's bit-length. None of
the three is supersingular; all are ordinary with prime group
order. The discriminating features:

\begin{itemize}
\item secp256k1 has $j = 0$ (CM by $\ZZ[\omega]$); P-256 and
      brainpoolP256r1 are non-CM (generic $j$).
\item secp256k1's coefficient $a = 0$; P-256's $a = -3$;
      brainpoolP256r1's $a$ is a specific 256-bit value.
\item secp256k1's group order has a near-power-of-2 form (refuted
      as the sole cause by the smaller Koblitz curves, but it may
      still contribute at 256 bits specifically).
\end{itemize}

The most natural follow-up: probe additional $j = 0$ curves at
256 bits to see whether the failure tracks $j = 0$. Candidates:
custom-constructed $j = 0$ Koblitz curves over different
near-256-bit primes.

\subsection{Question B: why does LLL uniformly fail at $\geq 384$ bits?}

The LLL iteration cap in our implementation is
$\text{cap} = 500 \cdot d^2 \cdot 8 + 10^4$ where $d$ is the
lattice dimension. For $d = 10$ (= $m + 2 = 8 + 2$), this gives
$\text{cap} \approx 4.1 \times 10^5$.

Classical LLL convergence is $O(d^4 \log B)$ where $B$ is the
max entry norm. With $d = 10$ and $B \approx n^2 \approx 2^{2 \cdot 384}
\approx 2^{768}$ for P-384, the expected iteration count is
$\sim 10000 \cdot 768 \approx 7.7 \times 10^6$ — \textbf{a factor
of 20 above the cap}. This strongly suggests the cap is simply
too small for $\geq 384$-bit curves.

The cleanest single test: raise the cap to $10^7$ and re-run on
P-384. If LLL converges, the cap was the cause. Then the cap
formula should be $O(d^2 \log B)$ rather than $O(d^2)$.

\subsection{Why secp256k1 alone fails at 256 bits (but not P-256)}

P-256's $n$ and secp256k1's $n$ are both 256-bit. Their basis
$B$-matrices have entries of comparable bit-length. So why does
LLL converge in $\sim$700 ms for P-256 but hit the cap for
secp256k1?

We do not have a closed-form answer. The candidate sources are:
the specific $(a_i, t_i)$ distribution induced by secp256k1's
trace $t \approx 2^{129}$; the
$\ZZ[\omega]$-decomposition of Frobenius $\pi = a + b\omega$;
the $r_i = (kG)_x \bmod n$ distribution being narrower for
secp256k1 (where $n$ is very close to $p$) than for P-256.

A theoretical Gram--Schmidt analysis is sketched in the companion
note \cite{ourcompanionauditgs} but does not yet pin down the
mechanism precisely. The next concrete step is to compute the
Gram-matrix eigenvalue spectrum empirically for both curves at
the same parameter set.

\section{Implications and outlook}

The reframed open question is sharper than the original. The
phenomenon's cause is somewhere in the intersection of:
specific arithmetic of secp256k1's $(p, n, t, a, b)$ tuple, the
specific cap parameter in the LLL implementation, and the basis-
construction step's sensitivity to coefficient size.

Importantly, the empirical observation \emph{remains valid} as a
practical attack-cost benchmark: any HNP attack reporting LLL costs
on a single curve cannot extrapolate to other curves without
revalidation. The refuted hypothesis only ruled out one
\emph{theoretical} explanation; it did not weaken the
\emph{empirical} claim.

\section{Implications}

\subsection{For practical ECDSA attacks}

The observed degeneracy is \emph{not} a security claim about
secp256k1: an attacker with biased ECDSA signatures and willing
to use BKZ-$\beta$ (with $\beta$ chosen large enough) will still
recover the key. LLL is the \emph{first-line} attack; the
degeneracy delays it but does not block more thorough lattice
reduction. The observed cost difference --- a factor of
$\sim 15\times$ between LLL convergence on P-256 and LLL failure on
secp256k1 --- does not translate into a security margin.

\subsection{For HNP-attack benchmarking literature}

We note that published HNP-on-ECDSA papers
(e.g.,~\cite{nguyen2002, mulder2013, aldaya2020}) often
demonstrate the attack on a single curve (often P-256 or
secp256k1) without head-to-head curve comparison. Our observation
suggests that LLL-convergence benchmarks are \emph{curve-specific}
in a way that has not been previously documented; reported
attack costs that generalise across curves should be treated with
caution.

\subsection{For the GLV-aware HNP programme}

This phenomenon affects the proposed GLV-aware Hidden-Number
Problem variant for secp256k1~\cite[\S 4]{ourcompanionnote}, which
plans to use a lattice of similar structure (with the addition of
half-scalar bookkeeping for the GLV decomposition). If the
underlying near-power-of-2 mechanism propagates, the GLV-aware
secp256k1 attack will need a different reduction strategy.

\section{Open questions}

\begin{enumerate}
\item \textbf{Theoretical}: derive the precise Gram--Schmidt
      eigenvalue structure of the Boneh--Venkatesan basis for $n$
      of near-power-of-2 form. Quantify the spectral gap (or its
      absence) that LLL's size-reduction inequality measures.
\item \textbf{Empirical}: replicate the head-to-head probe across
      all NIST curves (P-192, P-224, P-256, P-384, P-521), the
      brainpool family, and the Koblitz family (secp192k1,
      secp224k1, secp256k1). Tabulate convergence behaviour as a
      function of how close $n$ is to a power of 2.
\item \textbf{Algorithmic}: identify the smallest BKZ block size
      $\beta$ that converges where LLL fails. Cost growth from
      LLL to BKZ-$\beta$ as a function of $\beta$.
\item \textbf{Curve-design}: are there ECC curve standardisation
      criteria that should explicitly avoid near-power-of-2 group
      orders for HNP-robustness reasons? Currently no such
      criterion appears in published curve-selection rationales.
\end{enumerate}

\section{Conclusion}

We have documented a curve-specific LLL-convergence asymmetry
between secp256k1 and P-256 at standard HNP-attack parameters.
The phenomenon is empirically robust (3-of-3 vs 0-of-3 across
three independent seeds) and reproducible via the companion code.
It does not affect ECDSA security in practice (BKZ recovers
where LLL fails) but it does suggest that HNP benchmark costs
reported on one curve do not automatically transfer to another.
The underlying mechanism --- whether near-power-of-2 $n$ is
indeed the cause, or whether some deeper basis-structure feature
is at play --- is an open question for further work.

\begin{thebibliography}{99}

\bibitem{boneh1996} D.~Boneh, R.~Venkatesan, ``Hardness of computing
the most significant bits of secret keys in Diffie--Hellman and
related schemes'', CRYPTO 1996.

\bibitem{nguyen2002} P.~Nguyen, I.~Shparlinski, ``The insecurity of
the digital signature algorithm with partially known nonces'',
\emph{J. Cryptol.} 2002.

\bibitem{mulder2013} E.~De Mulder, M.~Hutter, M.~E.~Marson,
P.~Pearson, ``Using Bleichenbacher's solution to the hidden number
problem to attack nonce leaks in 384-bit ECDSA'', CHES 2013.

\bibitem{aldaya2020} A.~C.~Aldaya, B.~B.~Brumley et al, ``LadderLeak:
breaking ECDSA with less than one bit of nonce leakage'', CCS 2020.

\bibitem{fips186} NIST FIPS 186-4, ``Digital Signature Standard'',
2013.

\bibitem{sec2} Certicom Research, ``SEC 2: Recommended Elliptic Curve
Domain Parameters'', 2010.

\bibitem{ourcompanionnote} Anonymous companion note on a novel
GLV-aware Hidden-Number Problem for secp256k1.

\end{thebibliography}

\end{document}
